% Build from the repository root with: make demo
% Change theme=light to theme=dark to preview the other built-in theme.
\documentclass{article}
\usepackage[theme=light]{technica}
\usepackage{hyperref}

\hypersetup{
    pdftitle={Technica in practice},
    pdfauthor={Technica},
    pdfsubject={A demonstration of the Technica style system}
}
\technicadefinecallout{designrule}{Design rule}{TechnicaExplanation}

\title{Technica in practice}
\author{A short guide to reliable retries}
\date{Prose, code, and technical notes}

\begin{document}
\maketitle

Technical writing moves between explanation and implementation. This sample
uses a small retry policy to demonstrate Technica's typography, semantic prose,
code listings, terminal sessions, and callouts.

\section{Explain the idea}

A \term{retry budget} limits how often an operation may be attempted again.
The delay between attempts should grow, while the total amount of work remains
bounded. A simple policy starts at 100 milliseconds and doubles the delay up
to a fixed ceiling:
\[
    d_n = \min\bigl(100 \cdot 2^n,\;1600\bigr)\;\mathrm{ms},
    \qquad n = 0,1,2,\ldots
\]
Here, $n=0$ is the first retry. Five scheduled retries wait for a total of
$100 + 200 + 400 + 800 + 1600 = 3100$ milliseconds, excluding the time spent
performing the operation itself.

\subsection{Give names to the moving parts}

Store the policy in \filepath{config/retry.conf}. Use an API such as
\api{retry_delay} to calculate each delay, and inspect the sample with
\command{cat config/retry.conf}. In an editor, press \key{Ctrl} + \key{S}
to save a change.

\begin{note}
The initial request is not a retry. A budget of five retries permits up to
six attempts in total.
\end{note}

\begin{explanation}
Increasing the delay gives a busy service time to recover. The ceiling keeps
individual waits bounded; a separate budget bounds the number of attempts.
\end{explanation}

\newpage
\section{Show the implementation}
\label{sec:implementation}

The C++ example below calculates delays without sleeping or making network
requests. Applying the cap at each step keeps the arithmetic bounded.

The return type is \inlinecpp{std::chrono::milliseconds}.

\begin{cpp}[label=code:retry]{retry.cpp}
#include <algorithm>
#include <chrono>
#include <iostream>

using namespace std::chrono_literals;

std::chrono::milliseconds retry_delay(unsigned retry) {
    auto delay = 100ms;
    constexpr auto cap = 1600ms;
    while (retry > 0 && delay < cap) {
        delay = std::min(delay * 2, cap);
        --retry;
    }
    return delay;
}

int main() {
    for (unsigned n = 0; n < 5; ++n) {
        std::cout << retry_delay(n).count() << " ms\n";
    }
}
\end{cpp}

\begin{terminal}[title={Build and run}]
$ clang++ -std=c++17 retry.cpp -o retry
$ ./retry
100 ms
200 ms
400 ms
800 ms
1600 ms
\end{terminal}

\begin{warning}
A timeout does not prove that an operation failed. Retrying a payment or a
write can duplicate its effect unless the operation is safe to repeat.
\end{warning}

\newpage
\section{Connect the details}

Code block~\ref{code:retry} on page~\pageref{code:retry} implements the delay
rule. Keep the policy visible alongside the code so readers can compare the
configuration with its behavior.

\begin{codeblock}[label=code:config]{config/retry.conf}
initial_delay_ms = 100
maximum_delay_ms = 1600
maximum_retries = 5
\end{codeblock}

\subsection{Review the policy}

\begin{enumerate}
    \item Decide which failures are safe to retry.
    \item Set both an attempt limit and an overall time budget.
    \item Record attempts so repeated failures are observable.
\end{enumerate}

\begin{advice}
For many clients, add random variation to the delay. Otherwise, clients that
fail together may also retry together and create another burst of traffic.
\end{advice}

\begin{designrule}
Keep the timing calculation separate from the operation that waits. A pure
calculation is easier to inspect and test with exact expected values.
\end{designrule}

\subsection{Use the style in your own document}

The source of this sample uses the standard article class and the default page
dimensions. Start with the following preamble; choose the dark theme by
changing \inlinecode{light} to \inlinecode{dark}.

\begin{codeblock}{Minimal preamble}
\documentclass{article}
\usepackage[theme=light]{technica}
\usepackage{hyperref}
\end{codeblock}

The ordinary prose commands, code environments, and callouts all follow the
selected theme. Hyperlinks and references work when the document loads
\texttt{hyperref}; compile twice to resolve page and block references.

\end{document}
